% =====================================================================
%  CN2026Labels - Relatorio final (esqueleto)
%  Preencher cada seccao. Total entre 10 e 15 paginas.
% =====================================================================
\documentclass[11pt,a4paper]{article}

\usepackage[utf8]{inputenc}
\usepackage[T1]{fontenc}
\usepackage[portuguese]{babel}
\usepackage[a4paper,margin=2.3cm]{geometry}
\usepackage{microtype}
\usepackage{graphicx}
\usepackage{booktabs}
\usepackage{enumitem}
\usepackage{xcolor}
\usepackage{listings}
\usepackage{tikz}
\usepackage{hyperref}

\hypersetup{colorlinks=true,linkcolor=black,urlcolor=blue!60!black}

\title{\textbf{CN2026Labels} \\[2pt]
       \large Sistema el\'astico de dete\c{c}\~ao e tradu\c{c}\~ao de
       caracter\'isticas em imagens sobre Google Cloud Platform}
\author{Grupo G06 \\
        Instituto Superior de Engenharia de Lisboa \\
        LEIRT / LEIC / LEIM --- Computa\c{c}\~ao na Nuvem, ver\~ao 2025/2026}
\date{Maio de 2026}

\begin{document}
\maketitle

\begin{abstract}
% TODO: resumo de 5--10 linhas: o que o sistema faz, os servicos
%       GCP utilizados e as tres propriedades de alto nivel
%       (desacoplamento, elasticidade, simplicidade).
\end{abstract}

\tableofcontents
\vspace{1em}

% =====================================================================
\section{Introdu\c{c}\~ao e objetivos}
\label{sec:objetivos}
% TODO: reformular os objetivos do enunciado por palavras proprias.
%       Identificar os grupos SF e SG.

% =====================================================================
\section{Arquitetura}
\label{sec:arquitetura}
% TODO: incluir o diagrama de arquitetura (TikZ ou PNG).
\begin{figure}[h]\centering
% \includegraphics[width=0.9\linewidth]{arquitetura.png}
\caption{Arquitetura do sistema CN2026Labels (Figura~1 do enunciado).}
\label{fig:arch}
\end{figure}

\subsection{Componentes e responsabilidades}
% TODO: um paragrafo (ou linha de tabela) por componente:
%       Cloud Function de lookup, servidor gRPC (MIG), Pub/Sub,
%       Labels-App (MIG), Cloud Function de logging, Cloud Storage,
%       Firestore.

\subsection{Decis\~ao: dois servi\c{c}os gRPC no mesmo processo}
% TODO: separacao SF/SG, diferentes permissoes (IAM), um binario.

% =====================================================================
\section{Contrato protobuf}
\label{sec:contrato}
% TODO: incluir extrato do .proto e justificar:
%       client streaming na submissao, sub-mensagem LabelEntry,
%       int64 epoch ms em vez de Timestamp, Empty para GetStatus,
%       LabelResult unica para pendente e concluido.

% =====================================================================
\section{Fluxo de opera\c{c}\~oes}
\label{sec:fluxo}
% TODO: lista numerada (1, 2, 3.1, 3.2, 4.0, 4.1, 4.2, 5, 6, 7, 8)
%       espelhando a Figura~\ref{fig:arch}. Realcar o que e
%       sincrono e o que e assincrono.

% =====================================================================
\section{Formatos de dados e mensagens}
\label{sec:dados}

\subsection{Mensagem Pub/Sub}
% TODO: payload JSON + atributos.

\subsection{Documento Firestore}
% TODO: tabela com campo / tipo / escrito por / notas.

\subsection{Justifica\c{c}\~ao do \texttt{labelsIndex} desnormalizado}
% TODO: limite array_contains, sensibilidade a caixa, compromisso.

% =====================================================================
\section{Aspetos relevantes da implementa\c{c}\~ao}
\label{sec:impl}
% TODO: identificacao do pedido, stream da imagem, padrao work-queue,
%       reutilizacao dos clientes Google, configuracao por metadados,
%       elasticidade efetiva via ScaleService, lookup das instancias.

% =====================================================================
\section{Pressupostos assumidos}
\label{sec:pressupostos}
% TODO: papeis da conta de servico, gRPC em texto simples,
%       imagem apenas no Storage, requestId UUID, RUNNING-only
%       no lookup, semantica do zero em datas, filename so no
%       primeiro chunk.

% =====================================================================
\section{Compara\c{c}\~ao com decis\~oes alternativas}
\label{sec:alternativas}
% TODO: tabela com Decisao / Escolhida / Alternativa / Razao.

% =====================================================================
\section{Pontos de falha}
\label{sec:falhas}
% TODO: tabela cobrindo quota da Vision, falha do worker,
%       falha pontual do Translate, falha do blob, VM ainda nao a
%       aceitar ligacoes, entrega duplicada do Pub/Sub,
%       consistencia eventual do Firestore, esmagamento do MIG.

% =====================================================================
\section{S\'intese do procedimento de \emph{deployment}}
\label{sec:deploy}
% TODO: scripts 00..99 e uma linha por cada.

% =====================================================================
\section{Objetivos atingidos e n\~ao atingidos}
\label{sec:resultados}
% TODO: lista de pontos atingidos (parte A: submissao em stream,
%       consulta por id, consulta por label/data, download opcional;
%       parte B: scale-up/down dos dois MIG); lista do que foi
%       deixado de fora deliberadamente (TLS, autoscaling).

% =====================================================================
\section{Conclus\~ao}
\label{sec:conclusao}
% TODO: paragrafo final -- recapitulacao das escolhas que tornaram o
%       sistema elastico e resiliente, e os dois proximos passos
%       (TLS, autoscaling).

% =====================================================================
\begin{thebibliography}{9}
\bibitem{enunciado} \emph{Computa\c{c}\~ao na Nuvem -- Trabalho final}, ISEL, v1.0, 4 de maio de 2026.
\bibitem{vision}    \url{https://cloud.google.com/vision/docs/labels}
\bibitem{translate} \url{https://github.com/googleapis/google-cloud-java}
\bibitem{grpc}      \url{https://grpc.io/docs/languages/java/}
\bibitem{mig}       \url{https://cloud.google.com/compute/docs/instance-groups}
\end{thebibliography}

\end{document}
